% Gráfico 5. Hogares: renta disponible real por habitante y consumo privado (variación, %) y tasa de ahorro (% de la renta disponible), 2007-2025.
% Datos: ../datos/g05-ahorro-hogares.csv (Eurostat nasa_10_ki y nama_10_gdp). Capas: barras → cero → líneas → puntos.
\pgfplotstableread[col sep=comma,comment chars={\#}]{datos/g05-ahorro-hogares.csv}\gCinco
\begin{tikzpicture}[tcee]
  \begin{axis}[dicgraf,
      ybar, bar width=0.34cm,
      xmin=2006.4, xmax=2025.6, xtick={2007,2009,...,2025},
      ymin=-14, ymax=20, ytick={-12,-8,...,20},
      legend columns=2, dicapas]
    \capa{1}{\addplot[on layer=dicbarras, dicbarra=dicrenta, area legend] table[x=anio, y=renta_real_hab] {\gCinco}; \addlegendentry{Renta disponible real por habitante (\%)}}
    \capa{2}{\dicraya{(2006.4,0) (2025.6,0)}}
    \capa{3}{\addplot[sharp plot, diclinea=dicconspriv, leyendalinea] table[x=anio, y=consumo_real] {\gCinco}; \addlegendentry{Consumo privado real (\%)}}
    \capa{4}{\addplot[sharp plot, diclinea=dicahorro, leyendalinea] table[x=anio, y=tasa_ahorro] {\gCinco}; \addlegendentry{Tasa de ahorro (\% de la renta)}}
    \capa{5}{\addplot[on layer=dicpuntos, forget plot, sharp plot, dicpuntos=dicahorro] table[x=anio, y=tasa_ahorro] {\gCinco};}
  \end{axis}
\end{tikzpicture}
